\documentclass[11pt]{article}
\usepackage[margin=1in]{geometry}
\usepackage{amsmath,amssymb,amsthm}
\usepackage{graphicx}
\usepackage{hyperref}

\newtheorem{theorem}{Theorem}
\newtheorem{remark}{Remark}

\newcommand{\C}{\mathbb C}
\newcommand{\R}{\mathbb R}
\newcommand{\Bc}{\mathcal B}
\newcommand{\Ind}{\operatorname{Ind}}
\newcommand{\supp}{\operatorname{supp}}
\newcommand{\dbar}{\partial_{\bar z}}

\title{Exterior Winding-Zero Does Not Force Closedness for Differential Chains}
\author{Run \texttt{fa\_banach\_001}, lane 8}
\date{July 18, 2026}

\begin{document}
\maketitle

\section*{Source Problem}

The source is Harrison Pugh, \emph{Applications of Differential Chains to
Complex Analysis and Dynamics}, arXiv:1012.5542 \cite{PughThesis}.  Remark
3.7.2 follows the generalized Cauchy residue theorem and proposes that the
closedness assumption may be unnecessary:

\begin{center}
\includegraphics[width=\textwidth]{figures/open_problem_crop.png}
\end{center}

\noindent
The key sentence for this packet is the proposed implication that
\[
  \Ind_J(w)=0\quad(w\in U^c)
  \qquad\text{should imply}\qquad
  \partial J=0 .
\]
We give a compactly supported differential-chain counterexample to exactly
that implication.

\section*{Result}

\begin{theorem}\label{thm:main}
There are a bounded open set \(U\subset\C\) and a compactly supported
differential 1-chain \(J\in\hat{\Bc}_1^1(\R^2)\), with
\(\supp J\subset U\), such that
\[
  \Ind_J(w)=0\quad\text{for every }w\in U^c,
\]
but \(J\) is not closed:
\[
  \partial J\ne0.
\]
Consequently, the implication proposed in Remark 3.7.2 of the source paper is
false.
\end{theorem}

\section*{Proof Intuition}

The winding number only tests \(J\) against the Cauchy kernels
\((z-w)^{-1}dz\), and for \(w\notin U\) these are holomorphic near the support
of the chain.  A \(\dbar\)-exact density is invisible to all such holomorphic
tests by integration by parts.  But the same chain can still have a real
boundary, because closedness tests \(J\) against all exact real 1-forms, not
only holomorphic differentials.

\section*{Proof}

\begin{proof}
Let \(U\subset\C\) be a bounded open disk.  Choose a real-valued, nonconstant
bump function \(\phi\in C_c^\infty(U)\).  Define the smooth compactly supported
vector field
\[
  X=\frac12\left(\frac{\partial\phi}{\partial x}\partial_x
      +\frac{\partial\phi}{\partial y}\partial_y\right).
\]
Let \(J\) be the differential 1-chain represented by \(X\) against planar
Lebesgue measure \(dA\):
\[
  \int_J \omega=\int_U \iota_X\omega\,dA
\]
for every test 1-form \(\omega\).  This is the local compact-support version
of the source's vector-field construction of chains \(\xi_{X,\mu}\), which
places \(J\) in \(\hat{\Bc}_1^1\).

Write \(z=x+iy\).  Since \(dz(X)=\frac12(\phi_x+i\phi_y)=\dbar\phi\), for every
\(w\in U^c\) we have
\[
  \int_J \frac{dz}{z-w}
  =
  \int_U \frac{\dbar\phi(z)}{z-w}\,dA(z).
\]
The function \(z\mapsto (z-w)^{-1}\) is holomorphic on a neighborhood of
\(\supp\phi\).  Therefore
\[
  \frac{\dbar\phi(z)}{z-w}
  =
  \dbar\left(\frac{\phi(z)}{z-w}\right).
\]
Because \(\phi\) has compact support in \(U\), the last function is compactly
supported in the plane.  Hence
\[
  \int_U \dbar\left(\frac{\phi(z)}{z-w}\right)\,dA(z)=0,
\]
by applying Stokes' theorem, or simply by integrating the compactly supported
first partial derivatives over \(\R^2\).  It follows that
\[
  \Ind_J(w)=\frac{1}{2\pi i}\int_J\frac{dz}{z-w}=0
  \qquad(w\in U^c).
\]

Now test the boundary of \(J\) against the smooth 0-form \(\phi\).  Stokes'
formula for differential chains gives
\[
  \int_{\partial J}\phi=\int_J d\phi
  =\int_U d\phi(X)\,dA
  =\frac12\int_U |\nabla\phi|^2\,dA.
\]
Since \(\phi\) is not constant and has compact support, this integral is
strictly positive.  Therefore \(\partial J\ne0\).  This proves the claimed
counterexample.
\end{proof}

\section*{What This Does and Does Not Settle}

The counterexample disproves the specific mechanism proposed in Remark 3.7.2:
vanishing exterior winding number does not force a differential 1-chain to be
closed.  It does not, by itself, disprove the residue theorem after removing
closedness.  Also note that if the theorem's current assumption
\(\partial K=J\) is retained, then closedness follows automatically from
\(\partial^2K=0\).  Thus the remaining meaningful issue is not this implication
but whether the formula can be repaired or extended under a different
hypothesis, and whether the finite-mass requirement on \(K\) can be weakened.

\section*{Additional Diagnostic: Cauchy Formula Needs Closedness}

The thesis also states a generalized Cauchy integral formula without an
explicit closedness assumption on \(J\), and the later paper
\emph{Generalizations of the Cauchy Integral Theorems} repeats the same
formulation \cite{HarrisonPugh}.  The proof invokes the generalized Cauchy
integral theorem, which does require \(\partial J=0\).  A classical segment
already shows that some such assumption is necessary.

Indeed, let \(J\) be the oriented line segment from \(a\) to \(b\), with
\(a\ne b\), and let \(z\notin [a,b]\).  Put \(f(\zeta)=\zeta-z\).  Then
\(f(z)=0\), so the left side of the stated formula is \(0\).  But
\[
  \frac{1}{2\pi i}\int_J \frac{f(\zeta)}{\zeta-z}\,d\zeta
  =
  \frac{1}{2\pi i}\int_J 1\,d\zeta
  =
  \frac{b-a}{2\pi i}\ne0.
\]
This diagnostic is separate from Theorem~\ref{thm:main}; it is recorded here
because it points to the same missing closedness issue.

\section*{Novelty and Verification Notes}

The local run registry and cheap text indexes were searched for arXiv:1012.5542
and for the phrases ``Ind\_J(w)'', ``partial J'', ``closed'', and
``differential chains''.  The only local nearby result concerned a different
Banach flat-chain target, arXiv:2508.08017.  Web searches on 2026-07-18 for
the exact Remark 3.7.2 wording and close variants found the source thesis and
the later arXiv:1101.0385 paper, but no known counterexample to the implication
handled here.

For human review, the key technical point is membership of the vector-field
chain \(J=\xi_{X,dA}\) in the intended local differential-chain class.  The
source develops such chains from vector fields and Radon measures; with smooth
compact support this is the least singular case of that construction.

\begin{thebibliography}{9}

\bibitem{PughThesis}
Harrison Pugh,
\emph{Applications of Differential Chains to Complex Analysis and Dynamics},
arXiv:1012.5542, 2010.

\bibitem{HarrisonPugh}
Jenny Harrison and Harrison Pugh,
\emph{Generalizations of the Cauchy Integral Theorems},
arXiv:1101.0385, 2011.

\end{thebibliography}

\end{document}
